\chapter{Literature Review}

\section{Introduction}
This project sits between four bodies of work: continuous integration practice, automated program repair, the use of large language models for software engineering tasks, and the practical side of serving a model on modest hardware. This chapter summarises what each of them says that bears on our design, and ends with the gaps that the system tries to fill. Web documentation of the tools we used (GitHub Actions, n8n, llama.cpp and others) is cited where those tools are described in later chapters.

\section{Continuous Integration and Build Failures}
Continuous integration (CI) systems compile, build and test every change to a shared repository. Hilton et al.\ studied 34,544 open-source projects on GitHub and 1,529,291 builds, and surveyed 442 developers; their results support the claim that CI helps projects release more often and that its use keeps growing \cite{hilton16}. Shahin, Babar and Zhu reviewed the literature on continuous integration, delivery and deployment and collected its approaches, tools, challenges and practices \cite{shahin17}. Both works describe CI as an established practice, which is the setting of this project: a failing build is a routine event with a known cost in developer time, and a tool that shortens the time from a red build to a proposed fix addresses that cost directly.

Not every red build is a bug in the code. Luo et al.\ examined 201 commits that likely fix flaky tests in 51 open-source projects, where a flaky test is one whose outcome is not deterministic for the same code, and classified the most common root causes \cite{luo14}. For a repair tool this matters: asking a model to change the source in response to a flaky or infrastructure failure produces a patch for a problem that is not in the source. Our system acts on this by separating failures before any model call. The CI workflow marks a failure as \texttt{code}, \texttt{infra} or none, and the n8n healer refuses anything that is not a code failure (Chapter~5 and Chapter~6). The second healer, \texttt{ci-heal-backend}, goes further and classifies failures into four categories with a model before it decides between a retry, a fix and an issue (Chapter~8).

\section{Automated Program Repair}
Automated program repair (APR) is the study of producing patches without a person writing them. Gazzola, Micucci and Mariani surveyed 108 papers on automatic software repair and described the algorithms, compared them on representative examples and listed the open challenges \cite{gazzola17}. A classic technique is GenProg, which evolves a program variant using genetic programming and uses the existing test suite both to expose the defect and to protect required behaviour; delta debugging then reduces the difference to a minimal repair \cite{legoues12}. The idea that the test suite is the specification carries over unchanged into our design: a patch is accepted for review only if the pull request's own CI passes.

Xia, Wei and Zhang applied nine pre-trained language models, from 125M to 20B parameters, to repair tasks on five datasets in three languages and reported that directly applying the models could outperform existing APR techniques on all of those datasets \cite{xia23}. They also found that larger models tended to do better and that the code after the buggy line helps. This is the starting point for using a language model as the repair engine, and it explains why our prompt includes whole files and not only the failing line.

\section{LLM-Based Issue Resolution}
SWE-bench collects 2,294 software engineering problems drawn from real GitHub issues and their pull requests across 12 Python repositories, and asks a model to produce the patch that resolves each issue \cite{jimenez24}. It made repository-level repair a measurable task. SWE-agent showed that the interface given to a language-model agent matters: a purpose-built agent-computer interface improved the agent's ability to edit files, navigate a repository and run tests \cite{yang24}. Agentless then argued for a simpler design, a two-phase process of localisation followed by repair, without letting the model decide autonomously which tools to call \cite{xia25}.

Our pipeline is closest to the Agentless style. It has no tool-using agent loop. Localisation is done by ordinary code (the failing job's log and the files named in its stack trace or changed in the commit, capped at three candidates), the model is called once, and the repair is a small edit that code applies and checks. The retry loop adds a second pass when CI on the proposed fix still fails, but each pass is one prompt and one reply. The difference from the benchmarks above is the setting: the input is a CI log from a repository that already has a test suite, and the output is a pull request that CI evaluates, not a patch scored offline.

\section{Self-Healing and Autonomic Computing}
Kephart and Chess described autonomic computing as systems that manage themselves according to goals set by administrators, motivated by the difficulty of managing ever larger software environments \cite{kephart03}. A self-healing CI pipeline is a small instance of this idea: it monitors a signal (the build result), analyses it (is this a code failure, and what is the cause), plans an action (edit and open a pull request) and executes it, while a person keeps the final decision by merging. The capped retry loop in Chapter~6 can be read as the same monitor-analyse-plan-execute cycle repeated, with CI as the sensor.

\section{Serving a Large Model on Limited Hardware}
Running a model of about 27 billion parameters on two 16 GB GPUs relies on three techniques from the literature. \emph{Quantization} stores weights with fewer bits: Frantar et al.\ describe GPTQ, a one-shot method that reduces weights to 3 or 4 bits with small accuracy loss \cite{frantar22}. Our endpoint uses a roughly 4-bit GGUF file served by llama.cpp \cite{llamacpp}; the GGUF scheme is not GPTQ, but the motivation is the same. \emph{Speculative decoding}, introduced by Leviathan, Kalman and Matias, computes several tokens in parallel by letting a cheaper approximation propose them and the large model check them, without changing the output distribution \cite{leviathan23}. Gloeckle et al.\ trained models to predict several future tokens with separate output heads and reported that such models can decode up to three times faster \cite{gloeckle24}; the multi-token-prediction draft head used by our server follows this idea, and Chapter~7 reports its measured effect on our hardware (tokens per second with and without it). Kwon et al.\ showed how much serving throughput depends on key-value cache memory management \cite{kwon23}. Our server runs one request at a time, so its relevant limit is the cache size at a given context length, which Chapter~7 measures by stepping down from 131072 tokens.

\section{Reasoning Mode}
Wei et al.\ showed that prompting a model to produce intermediate reasoning steps improves its performance on arithmetic, commonsense and symbolic reasoning tasks \cite{wei22}. The Qwen3 model family integrates a thinking mode for multi-step reasoning and a non-thinking mode for fast answers into one model \cite{qwen3}. We use both: the first heal attempt runs with thinking off because it is fast, and the retry runs with thinking on because the first attempt failed. Chapter~9 measures what this choice costs in time and gains in fixes.

\section{Security of LLM-Integrated Applications}
Greshake et al.\ showed that applications which feed retrieved text into a model blur the line between data and instructions, so that an attacker who controls the retrieved text can steer the application \cite{greshake23}. The OWASP list of risks for LLM applications collects the same class of problems \cite{owasp-llm}. Our pipeline fits that description exactly: the prompt contains CI logs and source files that anyone who can push to the repository can influence. Chapter~10 describes the controls (the failure-kind gate, exact-match edits, a syntax check, CI on the pull request and human review) and states which attacks they do not stop.

\section{Tools and Platforms Used}
The pipeline is assembled from existing tools, and their documentation is the reference for the behaviour this report relies on. GitHub Actions defines workflows, jobs, runners and the failure conditions used in Chapter~5 \cite{gh-actions}; the GitHub REST API supplies the job list, the log download, the file contents and the branch and pull request calls \cite{gh-rest}; and the webhook documentation describes how a delivery is signed and sent \cite{gh-webhooks}. The healing workflow runs in n8n \cite{n8n-docs}. The model is served by the llama.cpp server \cite{llamacpp} on a Kaggle notebook \cite{kaggle-docs} and published through a Cloudflare quick tunnel \cite{cf-tunnel}. The second healer is written with FastAPI \cite{fastapi}, and the demo application that CI runs on uses Express \cite{express}.

\section{Summary of Research Gaps}
Table~\ref{tab:gaps} lists what the cited work leaves open and how this project responds.

\begin{table}[!htbp]
\centering
\small
\begin{tabularx}{\linewidth}{>{\raggedright\arraybackslash}p{4.2cm} >{\raggedright\arraybackslash}X >{\raggedright\arraybackslash}X}
\toprule
\textbf{Area} & \textbf{What the literature gives} & \textbf{What this project adds} \\
\midrule
CI failures & Usage data and flaky-test causes \cite{hilton16,luo14} & A gate that keeps infrastructure failures away from the model \\
Repair engines & Test-driven repair; LLM repair on benchmarks \cite{legoues12,xia23} & Repair triggered by a real CI log, judged by the repository's own CI \\
Issue resolution & Benchmarks and agent designs \cite{jimenez24,yang24,xia25} & A single-call design with code-side localisation and a capped retry loop \\
Serving & Quantization, speculative decoding, cache management \cite{frantar22,leviathan23,kwon23} & Measured speed and quality on two free 16 GB GPUs, with all flags and numbers recorded \\
Reasoning & Chain-of-thought and unified thinking models \cite{wei22,qwen3} & A measured policy: thinking off first, on for the retry \\
Security & Indirect prompt injection \cite{greshake23,owasp-llm} & A layered set of gates and an audit of the second healer's code \\
\bottomrule
\end{tabularx}
\caption{Gaps in the cited work and the response of this project.}
\label{tab:gaps}
\end{table}
